\documentclass{article}

\usepackage[margin=2.5cm]{geometry}
\usepackage{MiPsTeX}

\title{\textbf{MiPsTex: Example}} 
\author{Conde Lucento}
\date{2026}

\begin{document}
\maketitle
\tableofcontents
% ------------------------------
% Lista de códigos (Cod Index)
% ------------------------------
\listofcodes
\newpage

% ------------------------------
\section{Introducción}

\mipstitle{Objetivo del paquete MiPsTeX}

Este documento muestra todas las funcionalidades del paquete \texttt{MiPsTeX}:

\begin{itemize}
    \item Títulos y subtítulos estilo "pixel" (\verb|\mipstitle| y \verb|\submipstitle|)
    \item Entorno de código MIPS con estilo y colores
    \item Bloques de código opcionalmente titulados
    \item Índice automático de códigos (Cod Index)
\end{itemize}

\submipstitle{Primer bloque de ejemplo}

\begin{mips}[Inicialización de registros]
.text
main:
    li $t0, 10      # Cargar valor 10 en t0
    li $t1, 20      # Cargar valor 20 en t1
\end{mips}

% ------------------------------
\section{Operaciones básicas}

\mipstitle{Suma y Resta}

\submipstitle{Suma de dos números}

\begin{mips}[Suma de dos números]
.text
main:
    li $t0, 5
    li $t1, 7
    add $t2, $t0, $t1    # t2 = t0 + t1
    li $v0, 10
    syscall
\end{mips}

\submipstitle{Resta de dos números}

\begin{mips}[Resta de dos números]
.text
main:
    li $t0, 15
    li $t1, 4
    sub $t2, $t0, $t1    # t2 = t0 - t1
    li $v0, 10
    syscall
\end{mips}

% ------------------------------
\section{Acceso a memoria}

\mipstitle{Uso de memoria en MIPS}

\submipstitle{Cargar y almacenar palabras}

\begin{mips}[Acceso a memoria]
.data
num: .word 42

.text
main:
    la $t0, num        # Dirección de num en t0
    lw $t1, 0($t0)     # Cargar valor de num en t1
    li $v0, 10
    syscall
\end{mips}

\submipstitle{Almacenar valor en memoria}

\begin{mips}[Guardar en memoria]
.data
num: .word 0

.text
main:
    li $t0, 99
    la $t1, num
    sw $t0, 0($t1)    # Guardar t0 en la dirección de num
    li $v0, 10
    syscall
\end{mips}

% ------------------------------
\section{Bloques sin título}

\mipstitle{Bloques de código opcionales}

Este bloque de código no tiene título, por lo que **no aparecerá en el Cod Index**, pero mantiene los colores, fondo y numeración.

\begin{mips}
.text
main:
    li $a0, 100
    li $v0, 1
    syscall
    li $v0, 10
    syscall
\end{mips}

% ------------------------------
\section{Resumen de funcionalidades}

\submipstitle{Lista de características del paquete}

\begin{itemize}
    \item \textbf{Títulos grandes} con \verb|\mipstitle| y fondo negro
    \item \textbf{Subtítulos} con \verb|\submipstitle|, más pequeños
    \item \textbf{Entorno de código MIPS} con:
    \begin{itemize}
        \item Colores por tipo de instrucción
        \item Comentarios en verde
        \item Fondo gris
        \item Numeración de líneas
    \end{itemize}
    \item \textbf{Bloques titulados} agregan automáticamente entradas en el Cod Index
    \item \textbf{Bloques sin título} se formatean igual pero no aparecen en el índice
\end{itemize}

\end{document}